\insertImage{1d19cb7a-fad3-43c3-b5ed-42e53fc67784.png}{Aquarelle représentant diverses personnes dans un environnement professionnel dynamique et transparent. Leurs expressions sont satisfaites et épanouies, symbolisant leurs expériences positives.}

\chapter{Les expériences des employés}

\lettrine[lines=2]{J}{usqu'ici}, nous avons parlé structures, décisions, indicateurs. Il est temps de parler des seules personnes sans qui tout ça n'est que de la théorie d'organisation : celles qui vivent le modèle du lundi au vendredi.

Et je veux poser une règle pour ce chapitre : on écoute tout le monde. Pas seulement les enthousiastes qu'on fait monter sur scène dans les conférences, mais aussi ceux pour qui la libération a été une épreuve. Un livre qui s'appelle « l'envers du décor » et qui ne raconterait que des épanouissements ne vaudrait pas mieux que les plaquettes qu'il prétend dépasser.

\section{L'impact de l'autonomie sur l'engagement des employés}

Commençons par ce qui est solide : le lien entre autonomie et engagement n'est pas un slogan de consultant, c'est l'un des résultats les mieux établis de la psychologie du travail. Quand les gens contrôlent la manière dont ils font leur travail, leur motivation et leur implication augmentent — la recherche le mesure depuis des décennies\footnote{Langfred, C. W., \& Moye, N. A. (2004). Effects of task autonomy on performance: an extended model considering motivational, informational, and structural mechanisms. Journal of Applied Psychology, 89(6), 934-945.}.

Le terrain le confirme. Chez Chrono Flex, les techniciens itinérants ont reçu le pouvoir de décider sur place, chez le client, sans remonter la chaîne — et leur dirigeant décrit un effet immédiat sur la réactivité comme sur la fierté du métier\footnote{Gérard, A. (2017). Le patron qui ne voulait plus être chef. Flammarion.}. Rien de mystérieux là-dedans : être traité en adulte donne envie de se comporter en adulte. C'est toute la thèse du modèle, et sur ce point, elle tient.

Mais l'autonomie a une face B, et il faut la dire avec la même netteté : l'autonomie est une charge. Décider, c'est porter la décision. Là où l'organisation classique offrait le confort discret de la non-responsabilité — « je n'ai fait qu'appliquer » —, l'entreprise libérée expose chacun aux conséquences de ses choix. Pour certains, c'est grisant. Pour d'autres, c'est une source d'anxiété permanente, d'autant plus difficile à avouer que l'enthousiasme est devenu la norme sociale de l'entreprise. Quand « être autonome » est la valeur cardinale, dire « je suis perdu » devient un aveu coûteux.

Et il y a les préférences, tout simplement. Certains veulent un cadre, des attentes claires, une frontière nette entre le travail et la maison — et font par ailleurs un travail remarquable. Une entreprise libérée qui traite ces profils comme des salariés de seconde zone fabrique de la souffrance avec les meilleures intentions du monde. L'autonomie devrait être une offre. Malheur à ceux qui en font une injonction.

\section{Gestion des défis et des conflits dans un environnement de travail libéré}

Reprenons la question du conflit, mais cette fois du point de vue de celui qui le vit.

Vous êtes en désaccord profond avec un collègue — sur une décision, une répartition de la charge, une attitude. Dans une organisation classique, vous savez à qui en parler, même si vous savez aussi que ce sera imparfait. Dans une entreprise libérée mal outillée, vous êtes seul face à une question vertigineuse : à qui je m'adresse ? Il n'y a pas de chef. Aller voir l'équipe, c'est transformer un différend en tribunal. Ne rien dire, c'est laisser pourrir. La transparence, si fière d'elle-même, fait le reste : tout le monde voit tout, y compris vos tensions.

Ajoutez le problème de la décision évitée : quand personne n'a clairement l'autorité, certaines décisions inconfortables — recadrer un comportement, arrêter un projet auquel un collègue tient — ne sont prises par personne\footnote{Hollander, E. P. (1992). Leadership, followership, self, and others. Leadership Quarterly, 3(1), 43-54.}. Chacun attend que quelqu'un d'autre ose. Dans le vocabulaire des développeurs, c'est un deadlock : tout le monde attend tout le monde.

Les entreprises libérées matures répondent par des mécanismes explicites, et on retrouve ceux du chapitre sur les ressources humaines — médiation par les pairs, escalade progressive définie à l'avance, rôles de facilitation\footnote{Laloux, F. (2014). Reinventing Organizations: A Guide to Creating Organizations Inspired by the Next Stage of Human Consciousness. Nelson Parker.}. Je n'y reviens pas, sauf pour insister sur le point de vue du salarié : ce qui compte pour lui, ce n'est pas l'élégance du mécanisme, c'est de savoir \emph{avant d'en avoir besoin} qu'il existe, comment il marche, et qu'il peut s'en saisir sans être étiqueté fauteur de trouble. Un processus de médiation que personne n'ose déclencher est un poster de plus dans le hall.

\section{Ce que disent ceux qui y travaillent}

Je dois vous faire un aveu d'auteur : les « témoignages » qu'on lit dans la littérature sur les entreprises libérées — y compris dans la première version de ce chapitre — sont presque toujours de seconde main, souvent lissés, parfois inventés. J'ai décidé de ne pas jouer à ce jeu-là. Voici ce qu'on sait vraiment, et d'où on le sait.

Ce qu'on sait des enthousiasmes est bien documenté par ceux qui ont raconté leur transformation de l'intérieur. Les récits de FAVI rapportent des opérateurs fiers de gérer leur production et attachés à leur autonomie au point qu'un retour en arrière semblait impensable\footnote{Zobrist, J.-F. (2012). La belle histoire de FAVI: L'entreprise qui croit que l'homme est bon. Humanisme \& Organisations.}. Chez Chrono Flex, Alexandre Gérard décrit des équipes transformées par la confiance reçue — et ne cache d'ailleurs pas les turbulences du chemin\footnote{Gérard, A. (2017). Le patron qui ne voulait plus être chef. Flammarion.}.

Ce qu'on sait des difficultés est plus dispersé, parce que les déçus n'écrivent pas de livres : surcharge de réunions et de rôles chez les uns, anxiété de la responsabilité chez d'autres, sentiment d'un pouvoir devenu informel — donc inattaquable — chez d'autres encore. Ces récits-là, on les recueille en off.

Et puis il y a ce que j'ai entendu moi-même, dans les circonstances les plus honnêtes qui soient : la liquidation de ma propre société — j'y reviendrai dans la conclusion de ce livre. Quand une entreprise ferme, plus personne n'a de raison d'enjoliver quoi que ce soit ; ce qui se dit à ce moment-là, c'est la vérité résiduelle. Voici ce que mes collaborateurs m'ont dit en partant.

Plusieurs m'ont confié qu'ils allaient chercher ailleurs en sachant qu'ils ne retrouveraient jamais ce cadre-là — la pire des publicités pour le marché du travail, la meilleure pour le modèle. D'autres, plus jeunes, m'ont dit vouloir rejoindre des entreprises classiques, précisément pour comprendre ce qu'ils perdraient au change — et j'ai trouvé cette lucidité magnifique : on ne mesure une liberté qu'en la quittant. Et certains m'ont dit, sans détour, que trop de responsabilités, ce n'était pas pour eux.

Trois réactions, un seul modèle. Les trois sont vraies en même temps, et aucune n'annule les autres : c'est exactement ça, l'expérience des employés d'une entreprise libérée. Un salarié épanoui par l'autonomie et un salarié épuisé par elle peuvent travailler dans la même entreprise, la même année, au même étage. Tant qu'on n'accepte pas que ces deux vérités coexistent, on ne comprend rien à ce chapitre — ni à ce livre.
